% Part: computability
% Chapter: recursive-functions
% Section: pr-functions-computable

\documentclass[../../../include/open-logic-section]{subfiles}

\begin{document}

\olfileid[hi]{cmp}{rec}{cmp}
\olsection{आदिम पुनरावर्ती फलन संगणनीय हैं}

मान लें कि कोई फलन $h$ आदिम पुनरावर्तन द्वारा इस प्रकार परिभाषित है:
\begin{eqnarray*}
h(\vec x, 0)     & = & f(\vec x) \\
h(\vec x, y + 1) & = & g(\vec x, y, h(\vec x, y))
\end{eqnarray*}
और मान लें कि फलन $f$ तथा $g$ संगणनीय हैं। ($\vec x$ का उपयोग हम
$x_0$, \dots, $x_{k-1}$ को संक्षेप में लिखने के लिए करते हैं।)
तब $h(\vec
x, 0)$ को स्पष्टतः संगणित किया जा सकता है, क्योंकि वह केवल
$f(\vec x)$ है, जिसे हमने संगणनीय माना है। फिर $h(\vec x, 1)$ को भी
संगणित किया जा सकता है, क्योंकि $1 = 0 + 1$ और इसलिए $h(\vec x, 1)$
केवल यह है:
\begin{align*}
 h(\vec x, 1) & = g(\vec x, 0, h(\vec x, 0)) =  g(\vec x, 0, f(\vec x)).
\intertext{हम इसी प्रकार आगे बढ़कर इन्हें संगणित कर सकते हैं:}
h(\vec x, 2) & = g(\vec x, 1, h(\vec x, 1)) = g(\vec x, 1, g(\vec x, 0, f(\vec x)))\\
h(\vec x, 3) & = g(\vec x, 2, h(\vec x, 2)) = g(\vec x, 2, g(\vec x, 1, g(\vec x, 0, f(\vec x))))\\
h(\vec x, 4) & = g(\vec x, 3, h(\vec x, 3)) = g(\vec x, 3, g(\vec x, 2, g(\vec x, 1, g(\vec x, 0, f(\vec x)))))\\
& \vdots
\end{align*}
अतः सामान्यतः $h(\vec x, y)$ संगणित करने के लिए $h(\vec x, 0)$,
$h(\vec x, 1)$, \dots{} को क्रमशः संगणित करते जाएँ, जब तक हम
$h(\vec x, y)$ तक न पहुँचें।

इस प्रकार, यदि फलन $f$ और $g$ संगणनीय हैं, तो आदिम पुनरावर्ती परिभाषा
एक नया संगणनीय फलन देती है। इसी तरह, यदि फलन $f$ और $g_i$ संगणनीय हैं,
तो फलनों के संयोजन से भी संगणनीय फलन मिलता है।

चूँकि मूल फलन $\Zero$, $\Succ$ और $\Proj{n}{i}$ संगणनीय हैं, और संयोजन
तथा आदिम पुनरावर्तन संगणनीय फलनों से संगणनीय फलन देते हैं, इसलिए प्रत्येक
आदिम पुनरावर्ती फलन संगणनीय है।

\end{document}
